\centerline{\bf \Huge Тренировочная олимпиада}
\centerline{\bf \Huge 8 класс}

\begin{flushright}
        \scriptsize{Арбуз арбуз привет}
\end{flushright}

\problem{1}
Найдите сумму

$$
\frac{2^2}{1 \cdot 3}+\frac{4^2}{3 \cdot 5}+\frac{6^2}{5 \cdot 7}+\ldots+\frac{100^2}{99 \cdot 101}
$$


\problem{2}
Дан неравнобедренный треугольник $A B C$. На его сторонах $A B$ и $B C$ отмечены точки $K$ и $L$ такие, что прямые $K L$ и $A C$ параллельны. Отрезки $A L$ и $K C$ пересекаются в точке $S$. Известно, что $A K=A S$ и $K L=L C$. Докажите,что $A L=K B$.

\problem{3}
Медные шары массой $1$ г, $2$ г, $\ldots$, $9$ г выложили в ряд слева направо
по возрастанию массы. Затем один шар убрали. Как за два взвешивания
узнать, какого шара нет?

\problem{4}
Найдите сумму всех натуральных чисел $n$, для которых число $n^2+7 n+1$ является квадратом некоторого натурального числа.

\problem{5}
В ЭлМШ учится 40 красоток и 10 красавчиков. Каждый красавчик влюблен ровно в тех красоток, которые умнее его, либо ровно в тех, которые богаче его. Докажите, что у каких-то двух красоток одни и те же поклонники. (Про любых двух людей можно сказать, кто из них умнее. Если человек A умнее человека B и человек B умнее человека C, то верно, что A умнее C. Такие же свойства выполнены для богатства.)